% Source: https://artofproblemsolving.com/community/c3h3108511_joshys_1000th_post
% Stable identity: aops-topic-3108511-q04
If $ T $ is the number of minutes, then $\Pr(T>n)=1-(1-2^{-n})^2=2\cdot2^{-n}-4^{-n}$ for $ n\ge0 $. Summing the tail probabilities gives $ E(T)=2\cdot2-\dfrac{4}{3}=\dfrac{8}{3}$ minutes, or $\boxed{160\text{ seconds}}$. The tail-sum formula holds because an outcome of duration $ t $ contributes one to precisely the first $ t $ indicators $\mathbf1_{T>n}$.
